% V. Methodology. Every count, cap and hardware figure comes from the
% results files below through numbers.tex or a generated table.
% Feeds from: results/disk/corpus_profile.csv, results/web/dataset_profile.json,
% results/realio/device_*.json, results/gpu/gpu_env.json, and the
% provenance.json in each results directory.
% Table T1 (datasets) comes from paper/tables/make_tables.py.

\section{Methodology}\label{sec:method}

\todo{Table T1 (datasets): the oracleGeneral block traces in four corpora
with their reconstruction rules (MSR, Alibaba and Tencent records as
requests; FIU pages merged into requests with a sector cap), twin traces
and request caps, from results/disk/corpus\_profile.csv; and the
RetailRocket splits from results/web/dataset\_profile.json. Cite the trace
and dataset sources~\cite{narayanan2008write,koller2010iodedup,
li2020alibaba,zhang2020osca,snia_tencent_cbs,cmupdl_cachedatasets,
yang2020twemcache,libcachesim,retailrocket2017}.}

\todo{Protocols: chronological validation and test splits; paired and
cluster bootstrap CIs (method, resamples and seed from the PREREGISTERED.md
files); repeats per configuration; the training range of every trained
artifact, which is never scored on its own training segment.}

\todo{Hardware table: CPU, the HDD and SSD models, kernel, page-cache
reclaim policy, and the GPU with its contention record, from
results/realio/device\_*.json and results/gpu/gpu\_env.json. Software:
PyTorch~\cite{paszke2019pytorch} and
scikit-learn~\cite{pedregosa2011sklearn}.}
